\section{Purpose and theorem order}

The aim is to isolate a minimal mathematical architecture for guarded state change. The development order is
\begin{equation}
\boxed{
\begin{gathered}
\text{initial datum}\to
\text{invariant}\to
\text{validator}\to
\text{target support}\to
\text{curvature}\\
\to\text{proof admission}\to
\text{gate}\to
\text{fixed state}\to
\text{persistent record}\\
\to\text{positive-form completion}\to
\text{bilateral flow and observer jets}\\
\to\text{finite-to-infinite completion}\to
\text{minimal faithful observer}\to
\text{finite Hermitian certificate}
\end{gathered}}
\label{eq:chain}
\end{equation}
A later object is never used to redefine an earlier one after a defect has been observed. The finite record theorem remains self-contained. The operator extension is deliberately downstream: Hilbert space is obtained by completing a declared positive form, finite packets must converge on one recovered carrier, observer dimension is derived from target-relevant memory, and only then is a finite matrix allowed to decide the limiting sign.

\begin{longtable}{p{0.12\textwidth}p{0.35\textwidth}p{0.43\textwidth}}
\toprule
Result & Mathematical obligation & Consequence \\
\midrule
T1 & Structured initial datum & No accepted content is assumed. \\
T2 & Invariant transition bound & Local admissibility accumulates controllably. \\
T3 & Candidate--validator separation & A candidate cannot certify itself. \\
T4 & Exact target determination & Invisible target-changing directions are detected. \\
T5 & Exact completion and support & Missing target information is repaired explicitly and bounded. \\
T6 & Path-order curvature & Directional order leaves an exact commutator residue. \\
T7 & Finite proof admission & Licensed acyclic derivations close syntactically; sound rules lift closure semantically. \\
T8 & All-obligation threshold & Positive margin is necessary and sufficient; the midpoint threshold is robust. \\
T9 & Clock-independent fixed selector & Readiness is intrinsic to the state rather than periodic timing. \\
T10 & Disjoint status classification & Every typed event is exactly commit, hold, or reject. \\
T11 & Exact refinement & Valid correction transfers residue without changing the total object. \\
T12 & Replay-complete append-only persistence & Earlier states remain reconstructible from immutable journal prefixes. \\
T13 & Persistent-record theorem & The complete finite record follows from T1--T12. \\
T14 & Positive-form carrier completion & The Hilbert carrier is derived by quotient and completion rather than assumed without source structure. \\
T15 & Bilateral flow and observer jets & One generator produces the forward/backward flow, parity ladder, exact remainder, and nested recognition kernels. \\
T16 & Recognition-complete operator limit & Corrected finite packets converge with explicit Gram and signed-form error bounds. \\
T17 & No-hidden-memory completion & Vanishing faithfulness residual forces all limiting adverse memory into the finite observer range. \\
T18 & Variational minimal observer & The minimum faithful channel count is the target-relevant memory rank; Hilbert--Schmidt tails quantify under-observation. \\
T19 & Exact infinite-to-finite inertia & The finite Hermitian matrix carries the negative index, kernel, and normalized positivity margin of the limiting form. \\
T20 & Outward finite promotion & A finite matrix bound plus a proved completion error yields a lawful strict positivity certificate. \\
\bottomrule
\end{longtable}

\section{Structured initial datum}

\begin{definition}[Foundational datum]
Fix once and for all
\begin{equation}
\mathfrak F=(\mathfrak T,\mathfrak V,\Theta,\I_0,h_0),
\label{eq:foundation}
\end{equation}
where $\mathfrak T$ is a type system, $\mathfrak V$ a validator family, $\Theta$ a collection of gate parameters, $\I_0$ an invariant reference, and $h_0$ a lineage origin.
\end{definition}

\begin{definition}[Dynamic state record]
At causal index $k$, let
\begin{equation}
\Ldg_k=(\Wset_k,\E_k,\Cset_k,\Hset_k,\Jset_k,\Dset_k,\I_k,h_k),
\label{eq:record}
\end{equation}
where $\Wset_k$ is an immutable transition journal, $\E_k$ the set of addressable events, $\Cset_k$ the committed events, $\Hset_k$ the held events, $\Jset_k$ the rejected events, $\Dset_k$ a dependency relation, $\I_k$ the current invariant datum, and $h_k$ the current lineage identifier.
\end{definition}

\begin{definition}[Structured genesis pair]
The dynamic genesis state is
\begin{equation}
\Ldg_0=(\Wset_0,\E_0,\Cset_0,\Hset_0,\Jset_0,\Dset_0,\I_0,h_0),
\qquad
\Wset_0=\E_0=\Cset_0=\Hset_0=\Jset_0=\Dset_0=\varnothing.
\label{eq:initial}
\end{equation}
The complete genesis datum is the pair
\begin{equation}
\mathfrak G_0=(\mathfrak F,\Ldg_0).
\label{eq:genesis-pair}
\end{equation}
\end{definition}

\begin{theorem}[T1: non-vacuity and type consistency of genesis]
The genesis pair $\mathfrak G_0$ is mathematically nonempty even though no event has yet been committed, and $\Ldg_0$ has the same dynamic record type as every later $\Ldg_k$.
\end{theorem}

\begin{proof}
The fixed component $\mathfrak F$ contains $\mathfrak T$, $\mathfrak V$, $\Theta$, $\I_0$, and $h_0$, so the admissibility structure is nonempty. The dynamic component $\Ldg_0$ has exactly the fields displayed in \cref{eq:record}, with the event-bearing fields initialized to the empty set and the invariant and lineage fields initialized to $\I_0,h_0$. Thus no committed content is assumed, but the genesis state is a well-typed member of the same state family used at later indices.
\end{proof}

\section{Invariant-admissible transitions}

Let $(\X,d_\X)$ be a state space and $(\Z,d_\I)$ an invariant-value space. Let
\[
\I:\X\to\Z
\]
be fixed before evaluation. For an event $e$, let $F_e:\X\to\X$ be its candidate transition map.

\begin{definition}[Local invariant admissibility]
The transition $F_e$ is $\varepsilon_e$-admissible at $x$ when
\begin{equation}
d_\I\bigl(\I(F_ex),\I(x)\bigr)\le \varepsilon_e.
\label{eq:local-invariant}
\end{equation}
The exact case is $\varepsilon_e=0$.
\end{definition}

\begin{theorem}[T2: telescoping invariant bound]
Suppose $x_{k+1}=F_{e_k}x_k$ and every accepted transition satisfies \cref{eq:local-invariant}. Then for every $n\ge1$,
\begin{equation}
d_\I\bigl(\I(x_n),\I(x_0)\bigr)
\le \sum_{k=0}^{n-1}\varepsilon_{e_k}.
\label{eq:telescope}
\end{equation}
In particular, exact local invariance implies $\I(x_n)=\I(x_0)$ for all $n$.
\end{theorem}

\begin{proof}
Insert the intermediate values $\I(x_1),\ldots,\I(x_{n-1})$ and apply the triangle inequality. Each consecutive term is bounded by the corresponding $\varepsilon_{e_k}$. If all bounds vanish, every consecutive invariant value agrees.
\end{proof}
